% Ad Atticum 8.4 (ad-atticum-08-04) --- English translation
% Translated by Claude (Anthropic) from the Latin (Perseus).
% Style guide: STYLE.md.

\ciceroLetterOpener{CICERO ATTICO salutem}

\ciceroSection{1} That Dionysius of yours --- yours rather than ours, whose character I had got to know well enough, but still preferred to stand by your judgement of him rather than my own --- not even respecting the testimony you had so often given me on his behalf, has shown himself proud in a fortune that he supposed was about to be mine. Such fortune as I have, so far as it can be effected by human counsel, I will steer in its motion by a certain rational method. What honour from me, what attention has been wanting to him --- what commendation, even, to others, of a man generally despised? Sooner than not exalt him with my praises, I have preferred to have my own judgement reproved by my brother Quintus and commonly by everyone else; sooner than seek some other teacher for my Ciceros, I have undertaken to coach them under him myself. As for the letters --- ye immortal gods --- which I have sent to that man! How much honour, how much affection in them! You would say, by Hercules, that it was Dicaearchus, or Aristoxenus, that was being summoned, and not the single most garrulous human being alive and the least apt at teaching.

\ciceroSection{2} ``But he has a good memory.'' He will tell people that I have a better. To these letters of mine he sent back a response of a kind that I, for my part, have never sent to anyone whose case I would not take on. For always, if I can, if I am not blocked by a case undertaken earlier --- I have never to any defendant, however humble, however squalid, however guilty, however estranged from me, refused so abruptly as this man has, plainly and without any exception, cut me off. I have never met with anything more ungrateful; and in that vice no evil is missing. But enough, and too much, of this.

\ciceroSection{3} I have got the ship ready. I am still waiting nevertheless for a letter from you, to know what they have to answer to my deliberation. At Sulmo, Gaius Atius the Paelignian has opened the gates to Antony, although there were five cohorts there; Quintus Lucretius has fled from the place. You know: Gnaeus is making for Brundisium, the position abandoned. The business is finished.
